\begin{figure}[htbp]
	\centering
	\begin{tikzpicture}[>=Stealth, line width=0.8pt, scale=1.0]
		
		% ── Trave AB ──────────────────────────────────────────
		\draw[thick] (0,0) -- (9,0);
		\draw[dashed, thick] (0,-0.18) -- (9,-0.18);
		\fill (0,0) circle (2pt) node[above=2pt] {$A$};
		\fill (9,0) circle (2pt) node[above=2pt] {$B$};
		
		% ── Carico distribuito p(z): solo su [1,5] ────────────
		\draw[thick, domain=1:5, samples=60]
		plot (\x, {1.0 + 0.4*sin((\x-1)*50)});
		\foreach \x in {1.2, 1.7, 2.2, 2.7, 3.2,
			3.7, 4.2, 4.7}{
			\pgfmathsetmacro{\h}{1.0 + 0.4*sin((\x-1)*50)}
			\draw[->, thin] (\x, \h) -- (\x, 0);
		}
		\node[right=2pt] at (3.05, 1.6) {$p(z)$};
		
		% ── Carico distribuito q(z): solo su [5.5,9] ─────────
		\foreach \x in {5.7, 6.2, 6.7, 7.2, 7.7, 8.2, 8.7}{
			\draw[->, thin] (\x-0.35, 0.5) -- (\x, 0.5);
		}
		\draw[thick] (5.5, 0.5) -- (9, 0.5);
		\node[right=2pt] at (7.05, 0.75) {$q(z)$};
		
		% ── Carico distribuito c(z): solo su [1,5] ────────────
		\foreach \x in {1.3, 2.0, 2.7, 3.4, 4.1, 4.8}{
			\draw[->] (\x+0.18, 0.82)
			arc[start angle=0, end angle=180, radius=0.18];
		}
		\node[right=2pt] at (2.30, 0.65) {$c(z)$};
		
		% ── Forze concentrate in P_k ──────────────────────────
		\pgfmathsetmacro{\Pk}{5.5}
		\fill (\Pk, 0) circle (2.5pt);
		\node[above=0.5pt] at (\Pk, 0) {$P_k$};
		
		% Z(P_k): verso destra
		\draw[->, line width=1.4pt]
		(\Pk, 0.0) -- (\Pk+1.1, 0.0);
		\node[below=2pt] at (\Pk+0.55, 0) {$Z(P_k)$};
		
		% Y(P_k): verso il basso
		\draw[->, line width=1.4pt]
		(\Pk, 0) -- (\Pk, -1.0);
		\node[left=2pt] at (\Pk, -0.5) {$Y(P_k)$};
		
		% mu(P_k): coppia antioraria
		\draw[->] (\Pk+0.22, 0.55)
		arc[start angle=0, end angle=180, radius=0.22];
		\node[above=2pt] at (\Pk, 0.72) {$\mu(P_k)$};
		
		% ── Azioni in A ───────────────────────────────────────
		% Z_A: verso sinistra (reazione)
		\draw[<-, line width=1.4pt]
		(0, 0) -- (-1.1, 0);
		\node[above=2pt] at (-0.65, 0) {$Z_A$};
		
		% Y_A: verso il basso
		\draw[->, line width=1.4pt]
		(0, 0) -- (0, -1.0);
		\node[left=2pt] at (0, -0.5) {$Y_A$};
		
		% mu(A): coppia antioraria
		\draw[->] (0.22, 0.35)
		arc[start angle=0, end angle=180, radius=0.22];
		\node[above=2pt] at (0, 0.52) {$\mu(A)$};
		
		% ── Azioni in B ───────────────────────────────────────
		% Z_B: verso destra
		\draw[->, line width=1.4pt]
		(9, 0) -- (10.1, 0);
		\node[below=2pt] at (9.55, 0) {$Z_B$};
		
		% Y_B: verso il basso
		\draw[->, line width=1.4pt]
		(9, 0) -- (9, -1.0);
		\node[right=2pt] at (8.2, -0.5) {$Y_B$};
		
		% mu(B): coppia antioraria
		\draw[->] (9.22, 0.55)
		arc[start angle=0, end angle=180, radius=0.22];
		\node[above=2pt] at (9, 0.72) {$\mu(B)$};
		
	\end{tikzpicture}
	\caption{Azioni esterne su una trave rettilinea piana:
		carichi distribuiti $p(z)$, $q(z)$, $c(z)$ definiti
		a tratti; forze e coppia concentrate $Z(P_k)$, $Y(P_k)$,
		$\mu(P_k)$ in un punto interno $P_k$; forze e coppie
		agli estremi $Z_A$, $Y_A$, $\mu(A)$ in $A$ e
		$Z_B$, $Y_B$, $\mu(B)$ in $B$.}
	\label{fig:azioni_esterne}
\end{figure}
